\bi{\emph{Two-target amortization horizon.} A $2\times2$ access study compares no memory (N), program access (P), definition access (D), and both (B), transferring from user-badge to same-cutoff user-engagement and next-cutoff user-badge under a 12-call budget. P and D produce valid cross-task solutions, whereas N and B exhaust the budget. At the later cutoff all four succeed, with validation AUCs \MeANoneAUC/\MeAProgramAUC/\MeADefinitionAUC/\MeABothAUC. B uses \MeABothOnlineK\,k online tokens versus N's \MeANoneOnlineK\,k, but costs \MeABothTotalK\,k when source costs are spread over two targets. Consistent with the cost model, initial extraction costs require a longer reuse horizon for net savings.}
{\emph{两个目标的摊销范围。}一个$2\times2$访问实验比较无记忆、程序访问、定义访问及两者同时访问，从user-badge迁移到同截点user-engagement及下一截点user-badge，预算为12次调用。跨任务中P和D产生有效方案，N和B用完预算；下一截点四组均成功。B在线使用\MeABothOnlineK\,k token，N使用\MeANoneOnlineK\,k；将源成本分摊到两个目标后，B为\MeABothTotalK\,k。初始提取成本需要更长的复用范围才能实现净节省。}
